\documentclass[12pt,a4paper]{article}
\usepackage[utf8]{inputenc}
\usepackage[T1]{fontenc}
\usepackage[polish]{polski}
\usepackage{amsmath}
\usepackage{amsfonts}
\usepackage{amssymb}
\usepackage{booktabs}
\usepackage{geometry}
\usepackage{enumitem}
\usepackage{graphicx}
\usepackage{float}
\usepackage{listings}
\usepackage{xcolor}
\usepackage{tabularx}
\usepackage{array}
\usepackage{etoolbox}
\usepackage{url}

% lstlisting nigdy nie łamany między stronami
\BeforeBeginEnvironment{lstlisting}{\begin{minipage}{\linewidth}}
\AfterEndEnvironment{lstlisting}{\end{minipage}}

\geometry{margin=2cm}

\definecolor{codebg}{rgb}{0.97,0.97,0.97}
\definecolor{codekey}{rgb}{0.0,0.3,0.7}
\definecolor{codecomment}{rgb}{0.2,0.5,0.2}
\definecolor{codestring}{rgb}{0.6,0.1,0.1}

\lstset{
    language=Python,
    basicstyle=\ttfamily\small,
    keywordstyle=\color{codekey}\bfseries,
    commentstyle=\color{codecomment}\itshape,
    stringstyle=\color{codestring},
    backgroundcolor=\color{codebg},
    frame=single,
    framesep=4pt,
    breaklines=true,
    showstringspaces=false,
    numbers=left,
    numberstyle=\tiny\color{gray},
    xleftmargin=2em,
    tabsize=4,
    columns=fullflexible,
    literate=
        {ą}{{\k{a}}}1 {Ą}{{\k{A}}}1
        {ę}{{\k{e}}}1 {Ę}{{\k{E}}}1
        {ó}{{\'o}}1   {Ó}{{\'O}}1
        {ś}{{\'s}}1   {Ś}{{\'S}}1
        {ż}{{\.z}}1   {Ż}{{\.Z}}1
        {ź}{{\'z}}1   {Ź}{{\'Z}}1
        {ć}{{\'c}}1   {Ć}{{\'C}}1
        {ń}{{\'n}}1   {Ń}{{\'N}}1
        {ł}{{\l{}}}1  {Ł}{{\L{}}}1,
}

\title{Bezpiecze\'nstwo przesy{\l}ania danych, Zadanie~6}
\author{Piotr Gugnowski, wydzia{\l} EiTI, nr indeksu 320690}
\date{Czerwiec 2026}

\pagestyle{empty}

\begin{document}

\maketitle

% ===============================================================================
\section*{Opis sytuacji}
% ===============================================================================

Skaner zainstalowany przy jednostce master sieci Modbus RTU przechwycił
ramkę \texttt{FC~0x10} (Write Multiple Registers) z~nieczytelnym 128-bitowym
kodem dostępu. Znana jest jedynie suma kontrolna \textbf{CRC kodu~$= \texttt{0xA77C}$},
wyznaczona zgodnie ze standardową specyfikacją Modbus CRC-16.

\begin{table}[H]
\centering
\begin{tabular}{llll}
\toprule
\textbf{Pole} & \textbf{Bajty} & \textbf{Wartość} & \textbf{Opis} \\
\midrule
Adres slave              & 1 & \texttt{0x0A}   & adres urządzenia docelowego \\
Kod funkcji              & 1 & \texttt{0x10}   & Write Multiple Registers \\
Adres pierwszego rej.    & 2 & \texttt{0x0000} & rejestr startowy kodu \\
Adres ostatniego rej.    & 2 & \texttt{0x0008} & rejestry 0 do 7: kod, rejestr 8: CRC~kodu \\
Liczba bajtów danych     & 1 & \texttt{0x12}   & 18~B = 16~B kodu + 2~B CRC~kodu \\
\midrule
Rejestr 0 (kod)          & 2 & \texttt{0x????} & bajty kodu 1 i~2 \\
Rejestr 1 (kod)          & 2 & \texttt{0x????} & bajty kodu 3 i~4 \\
Rejestr 2 (kod)          & 2 & \texttt{0x????} & bajty kodu 5 i~6 \\
Rejestr 3 (kod)          & 2 & \texttt{0x????} & bajty kodu 7 i~8 \\
Rejestr 4 (kod)          & 2 & \texttt{0x????} & bajty kodu 9 i~10 \\
Rejestr 5 (kod)          & 2 & \texttt{0x????} & bajty kodu 11 i~12 \\
Rejestr 6 (kod)          & 2 & \texttt{0x????} & bajty kodu 13 i~14 \\
Rejestr 7 (kod)          & 2 & \texttt{0x????} & bajty kodu 15 i~16 \\
\midrule
Rejestr 8 (CRC kodu)     & 2 & \texttt{0xA77C} & CRC-16/Modbus z~16~B kodu \\
\midrule
CRC ramki Modbus         & 2 & \texttt{0x????} & CRC-16 całej ramki \\
\bottomrule
\end{tabular}
\caption{Przechwycona ramka Modbus RTU. Pole danych: 8 rejestrów kodu
         (16~B) + 1 rejestr CRC~kodu (2~B) = 18~B.}
\end{table}

% ===============================================================================
\section*{Zadanie 6a: Odtworzenie kodu dostępu (4~pkt)}
% ===============================================================================

\subsection*{Treść zadania}

Odtworzyć kod dostępu do zamka składu broni.

% -----------------------------------------------------------------------------
\subsection*{Analiza}

Jedyna znana informacja o~kodzie to jego 16-bitowa suma kontrolna
CRC-16/Modbus~$= \texttt{0xA77C}$.
Warunku $\mathrm{CRC16}(k) = \texttt{0xA77C}$ spełnia nie jeden, lecz
$2^{112}$ różnych 128-bitowych ciągów (dowód w~Zadaniu~6b).
Nie jest zatem możliwe \textbf{jednoznaczne} odtworzenie kodu, można
jedynie wygenerować \emph{jeden z~możliwych} kodów pasujących do danego CRC.

% -----------------------------------------------------------------------------
\subsection*{Metoda wyznaczenia przykładowego kodu}

Korzystamy z~bijekcyjności odwzorowania
$(b_{15},\, b_{16}) \mapsto \mathrm{CRC16}(\text{kod})$
przy ustalonych pierwszych 14~bajtach (wykazanej w~Zadaniu~6b).

\begin{enumerate}[label=\arabic*.]
  \item Wybieramy dowolnie \textbf{pierwsze} 14~bajtów kodu
        (bajty 1–14, lewa strona). Dla uproszczenia przyjmujemy zera:
        \[
          b_1 = b_2 = \cdots = b_{14} = \texttt{0x00}
        \]
        Dowolny inny wybór tych 14~bajtów również byłby poprawny
        i~dałby inną parę $(b_{15}, b_{16})$ spełniającą warunek CRC.
  \item Obliczamy stan rejestru CRC po przetworzeniu tych 14~bajtów.
        CRC jest algorytmem \textbf{sekwencyjnym}: stan po bajcie $i$ zależy
        wyłącznie od bajtów $b_1, \ldots, b_i$ i~nie wymaga znajomości
        kolejnych bajtów. Bajty $b_{15}$, $b_{16}$ jeszcze nie zostały
        przetworzone, więc stan $\mathrm{CRC\_state}_{14}$ liczymy
        niezależnie od nich
        (inicjalizacja \texttt{0xFFFF}, każdy bajt \texttt{0x00} daje XOR
        bez zmiany, jedynie iteracje LFSR modyfikują stan):

        \begin{center}
        \small
        \begin{tabular}{rll}
        \toprule
        \textbf{Nr bajtu} & \textbf{Bajt} & \textbf{Stan CRC po bajcie} \\
        \midrule
        init & & \texttt{0xFFFF} \\
        1  & \texttt{0x00} & \texttt{0x40BF} \\
        2  & \texttt{0x00} & \texttt{0xB001} \\
        3  & \texttt{0x00} & \texttt{0xC071} \\
        4  & \texttt{0x00} & \texttt{0x2400} \\
        5  & \texttt{0x00} & \texttt{0x0024} \\
        6  & \texttt{0x00} & \texttt{0x1B00} \\
        7  & \texttt{0x00} & \texttt{0x001B} \\
        8  & \texttt{0x00} & \texttt{0x0B40} \\
        9  & \texttt{0x00} & \texttt{0xF00A} \\
        10 & \texttt{0x00} & \texttt{0x0770} \\
        11 & \texttt{0x00} & \texttt{0xE406} \\
        12 & \texttt{0x00} & \texttt{0x0264} \\
        13 & \texttt{0x00} & \texttt{0xEB03} \\
        14 & \texttt{0x00} & \texttt{0x01AB} \\
        \bottomrule
        \end{tabular}
        \end{center}
        \[
          \mathrm{CRC\_state}_{14} = \texttt{0x01AB}
        \]
  \item Bajtów 15 i~16 \textbf{nie przepuszczamy przez LFSR jako niewiadomych},
        bo każda iteracja LFSR zależy od aktualnego LSB rejestru, który jest
        wyznaczany dopiero po XOR z~bajtem wejściowym. Zamiast tego
        \textbf{przeszukujemy} wszystkie $256 \times 256 = 65\,536$ możliwych
        kombinacji $(b_{15}, b_{16})$: dla każdej pary obliczamy CRC ze~stanu
        \texttt{0x01AB} i~sprawdzamy, czy wynik to \texttt{0xA77C}:
        \[
          \mathrm{crc16\_update}\bigl(\mathrm{crc16\_update}(\texttt{0x01AB},\,b_{15}),\,b_{16}\bigr)
          \stackrel{?}{=} \texttt{0xA77C}
        \]
  \item Wynik: $b_{15} = \texttt{0xF9}$, $b_{16} = \texttt{0x0A}$.
        Przebieg obliczeń:
        \[
          \texttt{0x01AB}
          \xrightarrow{\;b_{15}=\texttt{0xF9}\;}
          \texttt{0xFD80}
          \xrightarrow{\;b_{16}=\texttt{0x0A}\;}
          \texttt{0xA77C} \checkmark
        \]
\end{enumerate}

% -----------------------------------------------------------------------------
\subsection*{Przykładowy kod dostępu}

\begin{table}[H]
\centering
\begin{tabular}{lll}
\toprule
\textbf{Rejestr} & \textbf{Bajty (hex)} & \textbf{Opis} \\
\midrule
0 & \texttt{00 00} & bajty 1 i~2 \\
1 & \texttt{00 00} & bajty 3 i~4 \\
2 & \texttt{00 00} & bajty 5 i~6 \\
3 & \texttt{00 00} & bajty 7 i~8 \\
4 & \texttt{00 00} & bajty 9 i~10 \\
5 & \texttt{00 00} & bajty 11 i~12 \\
6 & \texttt{00 00} & bajty 13 i~14 \\
7 & \texttt{F9 0A} & bajty 15 i~16 (wyznaczone z~CRC) \\
\midrule
\multicolumn{2}{l}{Kod (hex): \texttt{0000000000000000000000000000F90A}} \\
\multicolumn{2}{l}{Weryfikacja: $\mathrm{CRC16}(\text{kod}) = \texttt{0xA77C}$} & $\checkmark$ \\
\bottomrule
\end{tabular}
\caption{Jeden z~$2^{112}$ możliwych kodów dostępu pasujących do
         CRC~$= \texttt{0xA77C}$.}
\end{table}

% -----------------------------------------------------------------------------
\subsection*{Kompletna odtworzona ramka}

Mając przykładowy kod, można uzupełnić całą ramkę Modbus.
CRC~ramki obliczane jest ze~wszystkich bajtów od~adresu slave do~CRC~kodu
włącznie (patrz Dodatek~B).

\begin{table}[H]
\centering
\small
\begin{tabular}{lll}
\toprule
\textbf{Pole} & \textbf{Bajty (hex)} & \textbf{Obliczone CRC po tym polu} \\
\midrule
Adres slave \texttt{0x0A}       & \texttt{0A}    & \texttt{0x473F} \\
Kod funkcji \texttt{0x10}       & \texttt{10}    & \texttt{0xDC06} \\
Adres reg. \texttt{0x0000}      & \texttt{00 00} & \texttt{0x3902} \\
Liczba reg. \texttt{0x0008}     & \texttt{00 08} & \texttt{0xB4C0} \\
Liczba bajtów \texttt{0x12}     & \texttt{12}    & \texttt{0x5D34} \\
Dane (14 bajtów \texttt{0x00})  & \texttt{00}$\times$14 & \texttt{0x228B} \\
Dane bajt 15: \texttt{0xF9}     & \texttt{F9}    & \texttt{0xC8A8} \\
Dane bajt 16: \texttt{0x0A}     & \texttt{0A}    & \texttt{0xB949} \\
CRC~kodu \texttt{0xA77C}        & \texttt{A7 7C} & \texttt{0xF38D} \\
\midrule
\textbf{CRC ramki}              & \textbf{\texttt{8D F3}} & (bajt niski, bajt wysoki) \\
\bottomrule
\end{tabular}
\caption{Krok po kroku obliczanie CRC ramki. Wartość CRC~$= \texttt{0xF38D}$;
         w~ramce: bajt niski \texttt{0x8D} jako pierwszy, bajt wysoki \texttt{0xF3}
         jako drugi.}
\end{table}

\noindent\textbf{Kompletna ramka (27 bajtów):}
\begin{center}
\ttfamily
0A~10~00~00~00~08~12~
00~00~00~00~00~00~00~00~00~00~00~00~00~00~F9~0A~
A7~7C~
8D~F3
\end{center}

% ===============================================================================
\section*{Zadanie 6b: Liczba możliwych kodów dostępu (4~pkt)}
% ===============================================================================

\subsection*{Treść zadania}

Ile jest możliwych kodów dostępu? Odpowiedź udowodnić.

% -----------------------------------------------------------------------------
\subsection*{Dowód}

\textbf{Twierdzenie.} Liczba 128-bitowych kodów dostępu $k \in \{0,1\}^{128}$
spełniających $\mathrm{CRC16}(k) = \texttt{0xA77C}$ wynosi dokładnie $2^{112}$.

\medskip
\noindent\textbf{Dowód.}

Podzielmy 16-bajtowy kod na dwie części:
\[
  k = (\underbrace{b_1,\ldots,b_{14}}_{\text{prefiks},\ 112\text{ bitów}},\;
       \underbrace{b_{15},\, b_{16}}_{\text{sufiks},\ 16\text{ bitów}})
\]

\noindent\textbf{Kluczowy lemat:} dla każdego wyboru prefiksu $(b_1,\ldots,b_{14})$
istnieje \textbf{dokładnie jeden} sufiks $(b_{15}, b_{16})$ taki, że
$\mathrm{CRC16}(k) = \texttt{0xA77C}$.

Jeśli lemat jest prawdziwy, zliczanie jest natychmiastowe:
\[
  \underbrace{2^{112}}_{\text{możliwych prefiksów}}
  \times
  \underbrace{1}_{\text{sufiks na prefiks}}
  = 2^{112}
\]

\noindent\textbf{Dowód lematu.}

CRC jest algorytmem sekwencyjnym: jedyną ``pamięcią'' po przetworzeniu prefiksu
jest bieżący stan rejestru $R \in \{0,\ldots,65535\}$ (liczba 16-bitowa).
Bez względu na to, jakie bajty tworzyły prefiks, dwa prefiksy prowadzące
do tego samego stanu $R$ są od tej chwili \textbf{nierozróżnialne} dla algorytmu.

Lemat jest zatem równoważny stwierdzeniu:
\begin{center}
\textit{Dla każdego stanu startowego $R$, odwzorowanie
$(b_{15}, b_{16}) \mapsto \mathrm{CRC16\_continue}(R,\, b_{15},\, b_{16})$
jest bijekcją (każdy wynik CRC pojawia się dokładnie raz).}
\end{center}

Stany startowe to liczby 16-bitowe, jest ich dokładnie $65\,536$.
Dla każdego z~nich sprawdzamy wszystkie $65\,536$ par $(b_{15}, b_{16})$
i~weryfikujemy, że wyniki są parami różne (każdy wynik jest jedną z~$65\,536$
możliwych wartości CRC, więc para różnych wejść dająca ten sam wynik oznaczałaby kolizję).
Łączna liczba obliczeń: $65\,536^2 \approx 4{,}3 \times 10^9$.
Każde obliczenie to dwa kroki LFSR z~tablicą lookup (jedno odczytanie z~pamięci
i~jeden XOR na krok), co umożliwia wykonanie całości w~czasie poniżej 4~sekund.
Jest to \textbf{skończone, wyczerpujące sprawdzenie} wykonane programem
\texttt{verify\_bijection.cpp} (kod źródłowy w~pliku \texttt{verify\_bijection.cpp},
wymagania i~instrukcja budowania w~pliku \texttt{README.md},
kompilacja na macOS: \texttt{just build}).

\medskip
\noindent Wynik działania programu \texttt{verify\_bijection}:

\begin{lstlisting}[language={}, numbers=none, basicstyle=\ttfamily\small]
Sprawdzono stanow: 65536
Stanow z kolizja:  0
Bijekcja dla wszystkich stanow: TAK
\end{lstlisting}

Dla każdego prefiksu istnieje zatem dokładnie jeden sufiks.
Prefiksów jest $2^{112}$, co kończy dowód:
\[
  \boxed{\text{Liczba możliwych kodów} = 2^{112}}
\]

% ===============================================================================
\section*{Zadanie 6c: Informacje do jednoznacznego wyznaczenia kodu (2~pkt)}
% ===============================================================================

\subsection*{Treść zadania}

Proszę o~komentarz, jakie informacje należałoby zdobyć, aby kod został
wyznaczony w~sposób jednoznaczny.

% -----------------------------------------------------------------------------
\subsection*{Odpowiedź}

Dysponujemy jedynie 16-bitowym CRC kodu. Kod ma 128 bitów, więc brakuje
$128 - 16 = 112$ bitów informacji do jednoznacznego wyznaczenia.
Aby zawęzić zbiór rozwiązań do jednego elementu, haker musiałby zdobyć
dodatkowe informacje eliminujące $2^{112} - 1$ fałszywych kandydatów.

\begin{enumerate}[label=\arabic*.]

  \item \textbf{Przechwycenie nieuszkodzonej ramki z~kodem.}
        Gdyby skaner zarejestrował ramkę bez zniekształceń, kod byłby wprost
        dostępny. Jest to najprostszy scenariusz: brakujące 128 bitów zostaje
        dostarczone bezpośrednio.

  \item \textbf{Przechwycenie dodatkowych bitów kodu z~innego źródła.}
        Dysponując co najmniej 112 dowolnymi bitami kodu (z~innej
        przechwytniętej ramki, przez dostęp fizyczny do urządzenia lub
        przez atak kanałem bocznym), można wyznaczyć pozostałe 16 bitów
        z~warunku CRC.

  \item \textbf{Wiedza o~strukturze lub formacie kodu.}
        Jeżeli kod nie jest losowy, lecz ma znany format (stały nagłówek,
        numer seryjny, ograniczony zakres wartości, generowany przez PRNG
        o~przestrzeni stanu $< 2^{112}$), każde takie ograniczenie zmniejsza
        liczbę kandydatów. Wystarczające ograniczenie prowadzi do
        jednoznaczności.

\end{enumerate}

\medskip
\noindent\textbf{Wniosek.}
Samo CRC-16 dostarcza tylko 16 bitów informacji o~128-bitowym kodzie.
Brakujące 112 bitów haker musi pozyskać z~innego źródła: bezpośrednio
(kolejna ramka, dostęp do urządzenia) lub pośrednio (wiedza o~strukturze kodu).


\end{document}
